\chapter{The sieve on words}\label{chap:words}

\jpoints{} finds the points over the triples $(a, b, c)$ of the main search
(Definition~\ref{def:mainTriples}) row by row: for each first coordinate $0 \le a \le H$ the
rows~$(a, b)$ of the box, those with $|b| \le H$, and $b \ge 1$ for $a = 0$, and in each row the third
coordinates $|c| \le H$. It sieves the third coordinates of a row as the bits of~$64$-bit words,
$c$ being bit $c \bmod 64$ of word $\lfloor c/64 \rfloor$ of the row. The sieving primes
$p_n = \mathrm{prime}[\mathrm{pnn}[n]]$, $n < s_3$, of the choice of the sieve
(Definition~\ref{def:IsChoice}; \lname{JPoints.sievingPrime}) come in three stages, $n < s_1$,
$s_1 \le n < s_2$ and $s_2 \le n < s_3$ (\texttt{sieve\_primes1}, \texttt{sieve\_primes2} and
\texttt{sieve\_primes3}):
\begin{enumerate}
  \item in the first stage the bit array of a chunk of the row is and-ed, word by word, with the
    words of the row of the class of~$(a, b)$ in the sieve table of each prime;
  \item in the second stage each word that the first one leaves nonzero is and-ed with the word of
    the row of each prime at the residue of the index of the word;
  \item in the third stage each bit that is left is tested on its own, by an entry of a table of
    the curve for each prime.
\end{enumerate}
A bit that survives is tested for $\gcd(a, b, c) = 1$ and goes to the exact check
(Lemma~\ref{lem:checkOnePoint_spec}). The bit array of a row starts from the mask of the row, which
holds the condition at~$2$, with the bits of $|c| > H$ at the ends of the row and the bit of the
third coordinate $c_0 = b^2/4a$ on the conic cleared: the triples of the conic $b^2 = 4ac$ are those
of the double points, which \texttt{find\_double\_points} searches on its own
(Section~\ref{sec:doublePoints}).

A component \emph{prints} a list of points when it appends their texts to the output
(Definition~\ref{def:Printf.emitPoint}), only the first under \texttt{-1}
(Definition~\ref{def:firstIf}); a component that returns a flag returns whether it stopped at a
point under \texttt{-1}. The sieve of a row of the box prints the points of the filtered enumeration, over the triples whose
third coordinate passes the test of the third coordinates (Definitions~\ref{def:colKeeps}
and~\ref{def:rowPoints}); that every point passes the tests is
Lemma~\ref{lem:colKeeps_of_hasPointOver}, from the classes of the points modulo the sieving primes
(Lemma~\ref{lem:CurveTables.admits12_of_hasPointOver}). The contracts are triples on the
components in \texttt{Id}, which take the state of the thread and return it, proved by
\texttt{vcgen}.

\section{The sieve tables of the first two stages}\label{sec:sieveTables}

For a sieving prime~$p$ of the first two stages, at most \texttt{MAX\_TABLE\_PRIME}, the classes
$(a, b, c)$ modulo~$p$ that the sieve admits (Definition~\ref{def:refAdmits12}) depend on~$c$ only
modulo~$p$. \texttt{init\_sieve} stores them in a table of~$p^2$ rows of~$p + 1$ words, the row of
the class~$(a, b)$, $a, b < p$, from word $(a p + b)(p + 1)$ on; bit~$j$ of word~$w$ of a row is
the class of $64 w + j$. As $64 p \equiv 0 \bmod p$, word~$i$ of a row of the search, which holds
the third coordinates $64 i + j$, has the bits of word $i \bmod p$ of the row of the table, and
word~$p$ repeats word~$0$.

The table of a prime is built in two steps (\lname{JPoints.initSievePrime}). The pattern of each
row, the bits of the classes $c < k p$ for the least multiple~$k p$ of~$p$ above~$64$, goes into a
row of the global array \texttt{help}, from the points at infinity, the values of~$f$ that are
squares and the exact table of the third stage (Lemma~\ref{lem:patternBits_iff}). Each pattern is
then unrolled over the~$p$~words of the row of the table: the words below $\lfloor k p/64 \rfloor$
are copied, the others come from a rotation of the words of the row of \texttt{help} as a ring, in
which the word at the cursor holds the $k p \bmod 64$ bits after the last window written and the
others the following~$64$~bits each (\lname{JPoints.initSieveRotate\_spec},
\lname{JPoints.initSieveRow\_spec}).

\begin{definition}[The sieve tables of the first two stages]\label{def:SieveTables.WF}
  \lean{JPoints.initSieve, JPoints.SieveTables.WF}
  \leanok
  \uses{def:CurveTables.WF}
  \texttt{init\_sieve} (\csource{j-points.c}{1087}{1089}) builds the table \texttt{sieve\_tab[n]}
  of each sieving prime $p = p_n$ of the first two stages, $n < s_2$, from the tables of
  \texttt{init\_main} and the tables~$ct$ of the curve. Sieve tables are \emph{well formed}
  for~$ct$, the indices \texttt{pnn} of the sieving primes in the table of the primes and~$s_2$,
  \lname{JPoints.SieveTables.WF}, when for every $n < s_2$, all $a, b < p$, $w \le p$ and $j < 64$,
  bit~$j$ of word~$w$ of the row of~$(a, b)$ in \texttt{sieve\_tab[n]}, the word
  $(a p + b)(p + 1) + w$, is the class test \lname{JPoints.CurveTables.admits12} of~$ct$ at the
  class of $(a, b, 64 w + j)$ modulo~$p$. In particular word~$p$ of a row repeats word~$0$.
\end{definition}

\formalnote{\texttt{help}, of static storage, which \texttt{init\_sieve} leaves rotated, is
threaded through the run with the tables. The sieve tables are stated by the class test read from
the tables of the curve, so that the contract of \texttt{init\_sieve} needs nothing of these
tables; that they are the references is Definition~\ref{def:CurveTables.WF}.}

\begin{lemma}[The bits of the sieve tables]\label{lem:SieveTables.WF.getLsbD_eq_admits12}
  \lean{JPoints.SieveTables.WF.getLsbD_eq_admits12}
  \leanok
  \uses{def:CurveTables.WF, def:SieveTables.WF}
  If the sieve tables are well formed for~$ct$, \texttt{pnn} and~$s_2$, $n < s_2$,
  $a, b < p = p_n$, $w \le p$ and $j < 64$, then bit~$j$ of word~$w$ of the row of~$(a, b)$ in the
  table of~$p$ is \lname{JPoints.CurveTables.admits12} of~$ct$ at the class of $(a, b, 64 w + j)$.
\end{lemma}

\begin{proof}
  \leanok
  It is the field of \lname{JPoints.SieveTables.WF}, through which the stages read the tables.
\end{proof}

\begin{lemma}[The patterns of the rows]\label{lem:patternBits_iff}
  \lean{JPoints.patternBits_iff}
  \leanok
  \uses{def:CurveTables.WF}
  Let $p = \mathrm{prime}[n]$ be a prime of the table, $a, b < p$, and~$k p$ the length of the
  patterns. The loops of \texttt{init\_sieve} set bit~$i$ of the pattern of the row~$(a, b)$ (all
  bits of the row~$(0, 0)$; for a curve with points at infinity modulo~$p$, the bits
  \lname{JPoints.InfinityBits} of the rows~$(0, b)$; the bits \lname{JPoints.OnJBits} of the
  exact table) exactly when $a = b = 0$, or $i < k p$ and the tables of the curve admit the class
  $(a, b, i)$ modulo~$p$ (\lname{JPoints.CurveTables.admits12}).
\end{lemma}

\begin{proof}
  \leanok
  \uses{def:refIsSquare}
  The loops set the bits of images: $i \equiv b c \bmod p$ in the row~$(0, b)$, $b \ne 0$,
  for~$f(c)$ a square, and $i \equiv a c \bmod p$ in the row $(a, a b \bmod p)$, $a \ne 0$, for a class
  $(1, b, c)$ of the exact table. Conversely a bit of a row~$(0, b)$ comes from $c = i b^{-1}$, and
  one of a row $(a, b')$ from the class $(1, b' a^{-1}, i a^{-1})$, with the inverses that
  \texttt{invert} computes, which are inverses by Fermat's little theorem
  (Definition~\ref{def:refIsSquare}).
\end{proof}

\begin{lemma}[The contract of \texttt{init\_sieve}]\label{lem:initSieve_spec}
  \lean{JPoints.initSieve_spec}
  \leanok
  \uses{def:Tables.WF, def:SieveTables.WF}
  If the tables of \texttt{init\_main} are well formed (Definition~\ref{def:Tables.WF}),
  $s_2 \le \texttt{NUM\_PRIMES}$ and the first~$s_2$ sieving primes are primes of the table of at
  most \texttt{MAX\_TABLE\_PRIME}, then \texttt{init\_sieve}, from the sieve tables of the start of
  the program, returns sieve tables that are well formed for the tables~$ct$ of the curve,
  \texttt{pnn} and~$s_2$, whatever~$ct$ is.
\end{lemma}

\begin{proof}
  \leanok
  \uses{def:refIsSquare, def:CurveTables.WF, lem:patternBits_iff}
  The loops write the patterns, whose bits are the classes admitted
  (Lemma~\ref{lem:patternBits_iff}), and the rows of the table are their windows
  (\lname{JPoints.initSievePrime\_spec}). The rows of \texttt{help} of distinct pairs of residues
  below \texttt{MAX\_TABLE\_PRIME} do not overlap, and the loop over the primes stores the tables
  in turn.
\end{proof}

\section{The static data, the state and the set-up}\label{sec:siftData}

\texttt{init\_sift} computes once for the run what the sieve of every row reads, its static data.
The thread keeps a state from row to row: the bit array, the rows of the sieve tables of the row
in hand with the statics of \texttt{row\_setup}, the data of the third stage and the state of the
analysis of the tube. The set-up of the search is what \texttt{find\_points} establishes after the
choice of the sieve, for every choice (Definition~\ref{def:IsChoice};
\lname{JPoints.findPointsSetup\_spec}); the components of the main search assume it.

\begin{definition}[The static data of the sieve]\label{def:SiftData.WF}
  \lean{JPoints.initSift, JPoints.SiftData.WF}
  \leanok
  \uses{def:CheckGlobals, def:Double.add, def:Double.ofDyadic}
  \texttt{init\_sift} (\csource{j-sift.c}{132}{141}) computes the static data of the sieve once
  for the run, after the choice of the primes and \texttt{init\_sieve}. Static data are those of
  the globals~$g$ of the sieve, \lname{JPoints.SiftData.WF}, when they hold:
  \begin{enumerate}
    \item the~$s_3$ sieving primes~$p_n$, the lengths $p_n + 1$ of the rows of their tables and
      $p_n (p_n + 1)$ of the blocks of the rows of one class of~$a$;
    \item the words of a whole row, from $w_{\mathrm{low}} = \lfloor -H/64 \rfloor$ to
      $w_{\mathrm{high}} = \lfloor H/64 \rfloor + 1$, exclusive, and the
      $\lceil (w_{\mathrm{high}} - w_{\mathrm{low}})/\texttt{array\_size} \rceil$ chunks of a row,
      the~$k$-th from word $w_{\mathrm{low}} + k\,\texttt{array\_size}$ to the least of
      $w_{\mathrm{low}} + (k + 1)\,\texttt{array\_size}$ and~$w_{\mathrm{high}}$, each with the
      walk of every first-stage prime over it (\lname{JPoints.initSiftWalk}) and the residues of
      its first word modulo the sieving primes;
    \item the residues $k \bmod p_n$ of the positions $k \le \texttt{array\_size}$ in a chunk;
    \item the row \texttt{all\_ones} of the third stage, which admits every class, and the
      reciprocals $\lceil 2^{32}/p \rceil$ of the third-stage primes;
    \item the bound on the fourth coordinate as a double and its square, and the
      reciprocals~$1/p_n$ as doubles.
  \end{enumerate}
\end{definition}

\begin{lemma}[The contract of \texttt{init\_sift}]\label{lem:initSift_spec}
  \lean{JPoints.initSift_spec}
  \leanok
  \uses{def:SiftData.WF}
  If the arrays of the static data have the sizes that \texttt{init\_sift} writes into
  (\lname{JPoints.SiftData.Sized}), $s_1 \le s_2 \le s_3 \le \texttt{NUM\_PRIMES}$ and the sieving
  primes are primes of the table, then \texttt{init\_sift} returns static data of the globals~$g$
  of the sieve.
\end{lemma}

\begin{proof}
  \leanok
  \uses{def:CheckGlobals, def:Double.add, def:Double.ofDyadic}
  From the contracts of its parts: the data of the sieving primes, the chunks with their walks,
  stated by the value of \lname{JPoints.initSiftWalk}, and the residues of the positions with the
  reciprocals.
\end{proof}

\begin{definition}[The state of the thread]\label{def:SiftState.WF}
  \lean{JPoints.SiftState.WF}
  \leanok
  A state of the thread (\lname{JPoints.SiftState}) is consistent for the globals of the sieve by
  bands and the static data, \lname{JPoints.SiftState.WF}, when its bit array has
  $\texttt{array\_size} + 1$ words, the statics of \texttt{row\_setup} are consistent with the rows
  of the sieve tables that it holds (\lname{JPoints.RowSetupCache.WF}), the data of the third stage
  are consistent (\lname{JPoints.Row3Cache.WF}), the state of the analysis of the tube is well
  formed, and the chunk looked up last is a chunk of a row. The counters and the output are not
  constrained.
\end{definition}

\formalnote{No contract says anything of the counters of the survivors of the three stages, which
go into the totals of the report.}

\begin{definition}[The set-up of the search]\label{def:IsSearchSetup}
  \lean{JPoints.IsSearchSetup, JPoints.IsSieveSetup}
  \leanok
  \uses{def:CheckGlobals, def:CurveTables.WF, def:TwoAdic.Sound, def:SieveTables.WF, def:SiftData.WF, def:IsSearchGlobals}
  What \texttt{find\_points} sets up after the choice of the sieve, for every value of the
  \texttt{Float}s of the cost model, \lname{JPoints.IsSearchSetup} of the globals~$g$ of the
  search, the globals of the sieve by bands and the static data:
  \begin{enumerate}
    \item stages of distinct sieving primes among the primes considered,
      $s_1 \le s_2 \le s_3 \le \texttt{num\_primes}$, those of the first two stages at most
      \texttt{MAX\_TABLE\_PRIME}, and none when the value of \texttt{-N} is at most~$0$;
    \item the globals of the exact check (Definition~\ref{def:CheckGlobals}) and the tables of
      \texttt{init\_main} of~$g$, and the tables of the curve for the primes considered and the
      sieving primes (Definition~\ref{def:CurveTables.WF});
    \item the mode~$0$ (the box), $1$ (the tube) or~$2$ (the real region), the tube only for a
      bounded fourth coordinate, and the box without the condition at~$2$ when there are no
      sieving primes;
    \item well-formed sieve tables (Definition~\ref{def:SieveTables.WF}), the static data of the
      sieve (Definition~\ref{def:SiftData.WF}) and a condition at~$2$ that admits every point
      (Definition~\ref{def:TwoAdic.Sound});
    \item the intervals on which~$f$ is negative, the coefficients of the tube, the size of the
      bit array and the masks of the first and the last word of a row of~$g$.
  \end{enumerate}
  \lname{JPoints.IsSieveSetup} of the globals of the sieve by bands and the static data says that
  there are globals of the search as \texttt{main} sets them (Definition~\ref{def:IsSearchGlobals})
  for which \lname{JPoints.IsSearchSetup} holds: what the components of the main search, which do
  not see the globals of the search, may assume.
\end{definition}

\begin{definition}[The set-up of the stages and of a row]\label{def:IsStageSetup}
  \lean{JPoints.IsStageSetup, JPoints.IsRowSetup}
  \leanok
  \uses{def:CheckGlobals.WF, def:Tables.WF, def:SieveTables.WF, def:SiftData.WF}
  What the stages read of the globals and the static data of the sieve,
  \lname{JPoints.IsStageSetup}: the globals of the exact check as \texttt{main} sets them
  (Definition~\ref{def:CheckGlobals.WF}), the tables of \texttt{init\_main}
  (Definition~\ref{def:Tables.WF}), well-formed sieve tables of the first two stages, the static
  data of the globals, and $s_1 \le s_2 \le s_3 \le \texttt{NUM\_PRIMES}$ with sieving primes of
  the table. What the sieve of a row reads, \lname{JPoints.IsRowSetup}, adds a bit array of at
  least one word and the masks \texttt{begmask} and \texttt{endmask} of the first and the last
  word of a row as \texttt{main} sets them.
\end{definition}

\section{The sieve of a row}\label{sec:siftRow}

For a row~$(a, b)$ of the box, \texttt{sift} first computes the third coordinate~$c_0$ of the conic
and the mask of the row (\lname{JPoints.siftMask}): $c_0 = b^2/4a$ when $a \ne 0$ and
$4 a \mid b^2$, and otherwise~$H + 64$, whose word lies beyond the row; the mask is the word of the
condition at~$2$ for $(a \bmod 64, b \bmod 64)$ when the sieve uses it and $a \ne 0$, else the
odd~$c$ when~$a$ and~$b$ are even, else all~$c$. Then it moves the rows of the sieve tables
to~$(a, b)$ and sieves each chunk of the row. \texttt{sift\_range} sieves a range of words of a chunk,
for the sieve by bands of the tube and the real region (Definition~\ref{def:siftBands}).

\begin{definition}[The sieve of a row]\label{def:sift}
  \lean{JPoints.rowSetup, JPoints.scanWords, JPoints.checkPoint, JPoints.sift0, JPoints.sift, JPoints.siftRange}
  \leanok
  \uses{def:CheckGlobals, def:Double.add, def:Double.ofDyadic, def:checkOnePoint}
  The components of the sieve of a row~$(a, b)$:
  \begin{enumerate}
    \item \lname{JPoints.rowSetup}, \texttt{row\_setup} (\csource{j-sift.c}{1109}{1110}): the rows
      of the sieve tables of the first two stages for~$(a, b)$;
    \item \lname{JPoints.scanWords}, \texttt{scan\_words} (\csource{j-sift.c}{959}{969}): each
      nonzero word of a range of positions of a chunk handed to \lname{JPoints.checkPoint},
      \texttt{check\_point} (\csource{j-sift.c}{357}{370}), the second stage of a word, whose bits
      left go to the third stage, the test of coprimality and the exact check;
    \item \lname{JPoints.sift0}, \texttt{sift0} (\csource{j-sift.c}{997}{1001}): the sieve of
      the~$k$-th chunk of the row, whole: the bit array filled with the mask, and-ed in the same pass
      with the pattern of the first first-stage prime when there is one, the bits outside the row and the bit of~$c_0$ cleared,
      the patterns of the other first-stage primes and-ed in by their walks, and the words handed
      to \texttt{scan\_words};
    \item \lname{JPoints.sift}, \texttt{sift} (\csource{j-sift.c}{1156}{1160}): the sieve of a row
      of the box, \texttt{row\_setup} and then every chunk by \texttt{sift0}, until a point is
      found under \texttt{-1};
    \item \lname{JPoints.siftRange}, \texttt{sift\_range} (\csource{j-sift.c}{1049}{1054}): the
      sieve of the words $[lo, hi)$ of the~$k$-th chunk, with walks that start inside the chunk.
  \end{enumerate}
\end{definition}

\begin{definition}[The points of a range of words]\label{def:rowRangePoints}
  \lean{JPoints.rowRangePoints}
  \leanok
  \uses{def:orderedTriples, def:HasPointOver, def:colKeeps, def:CheckGlobals}
  For the globals~$g$ and the static data~$sd$ of the sieve, a row~$(a, b)$ and words~$lo$
  and~$hi$, the list $\mathrm{rowRangePoints}(g, sd, a, b, lo, hi)$ holds the points that the exact
  check finds over the triples $(a, b, c)$ (Definition~\ref{def:CheckGlobals}), $|c| \le H$,
  $c$ increasing, whose third coordinate passes the test of the third coordinates
  \lname{JPoints.colKeeps} (Definition~\ref{def:colKeeps}), lies off the conic, $b^2 \ne 4 a c$,
  and has its word $\lfloor c/64 \rfloor$ in $[lo, hi)$.
\end{definition}

Over all the words of a row these are the points of the filtered enumeration
(Definition~\ref{def:rowPoints}; \lname{JPoints.rowPoints\_eq\_rowRangePoints}), and adjacent ranges
concatenate: the sieve of a row may proceed chunk by chunk and range by range.

\begin{lemma}[The contract of \texttt{row\_setup}]\label{lem:rowSetup_spec}
  \lean{JPoints.rowSetup_spec}
  \leanok
  \uses{def:IsStageSetup, def:sift}
  Let the set-up of the stages hold (Definition~\ref{def:IsStageSetup}), the statics of
  \texttt{row\_setup} in the state be consistent with its rows, $a \ge 0$ and $|b| \le H$. Then
  \texttt{row\_setup} returns the state with only its rows and its statics changed, these
  consistent for the last row~$(a, b)$ (\lname{JPoints.RowSetupCache.WF}): for each $n < s_2$, with
  $p = p_n$, the residue $b \bmod p$ and the index $((a \bmod p)\,p + (b \bmod p))(p + 1)$ of the
  row of~$(a, b)$ in the table of~$p$. For the last row $(a, b')$ it moves them on by the small
  step for $0 < b - b' < 3$ and by the residues of the step $d = b - b'$ for a larger step;
  otherwise it computes them by division.
\end{lemma}

\begin{proof}
  \leanok
  \uses{def:CheckGlobals, def:CheckGlobals.WF, def:SiftData.WF, def:Double.toRat, def:Double.ofDyadic, lem:Double.ofDyadic_exact, def:Double.add, def:Double.IsRounding, lem:Double.isRounding_ofDyadic, lem:Double.isRounding_add}
  A step between two rows of the box is below~$2^{48}$, as
  $|b| \le H \le \texttt{MAX\_HEIGHT} < 2^{32}$ (Definition~\ref{def:CheckGlobals.WF}) and the
  statics keep the last second coordinate within the height bound. The residues $d \bmod p$, which
  \texttt{row\_setup} computes once for each step other than the last, come from
  the product of~$d$ with the double~$1/p$: $d$ converts exactly
  (Lemma~\ref{lem:Double.ofDyadic_exact}), the reciprocal and the product are rounded to nearest
  (Lemma~\ref{lem:Double.isRounding_add}, for the quotient and the product), and the product
  truncates to a quotient~$q$ with $d - p q \in [d \bmod p, d \bmod p + p]$, which one correction brings to
  $d \bmod p$.
\end{proof}

\section{The second and the third stage}\label{sec:stages23}

A word that the first stage leaves nonzero goes to \texttt{check\_point} with its index~$i$ in the
row and its position~$k$ in the chunk. For each second-stage prime~$p$ the residue $i \bmod p$ is
the sum of $k \bmod p$ and of the residue of the first word of the chunk, after one conditional
subtraction (\lname{JPoints.WordResidues}). \texttt{check\_bits} (\csource{j-sift.c}{274}{290})
takes the bits that are left from the lowest, $c = 64 i + j$, and tests $\gcd(a, b, c) = 1$ by
\texttt{relprime3}, then the third stage by \texttt{check\_bit}, before the exact check.

The third stage tests each third-stage prime~$p$ by one entry of a table of the curve
(\csource{j-sift.c}{254}{267}). For the row~$(a, b)$, \texttt{row3\_setup} chooses a factor and a
row: for $p \nmid a$ the factor~$a^{-1}$ and the row of the class $(1, b a^{-1})$ of the exact
table; for $p \mid a$ and $p \nmid b$ the factor~$b^{-1}$ and the table of the values of~$f$ that
are squares; for~$p$ dividing both the row \texttt{all\_ones}. \texttt{check\_bit} reads the row at
the class of~$c$ times the factor, reduced modulo~$p$ by the reciprocal $\lceil 2^{32}/p \rceil$.
This is the test of the first two stages without the test of the points at infinity
(\lname{JPoints.CurveTables.admits3}), which admits every class that the first two admit.

\begin{lemma}[The test of coprimality]\label{lem:relprime3_spec}
  \lean{JPoints.oddGcd_spec, JPoints.relprime3_spec}
  \leanok
  \uses{def:colKeeps}
  For odd $n, m > 0$ the binary algorithm \lname{JPoints.oddGcd}, which replaces the larger of the
  two by the odd part of their difference while they differ, returns $\gcd(n, m)$. For $a \ge 0$,
  \texttt{relprime3} returns whether $\gcd(a, b, c) = 1$ (\lname{JPoints.coprimeKeeps},
  Definition~\ref{def:colKeeps}).
\end{lemma}

\begin{proof}
  \leanok
  The loop keeps the gcd and the oddness, and~$n + m$ decreases. \texttt{relprime3} answers false
  when the three coordinates are even, uses \texttt{relprime2} when one of them is~$0$, and
  otherwise takes the gcd of their odd parts, which is that of the triple when one of them is odd
  (\lname{Int.gcd\_gcd\_eq\_of\_two\_pow\_mul\_eq}).
\end{proof}

\begin{lemma}[The contract of \texttt{check\_bit}]\label{lem:checkBit_spec}
  \lean{JPoints.checkBit_spec}
  \leanok
  \uses{def:stage1Keeps, def:IsStageSetup}
  Let the set-up of the stages hold, the data of the third stage be consistent
  (\lname{JPoints.Row3Cache.WF}), the residues be those of the index $\lfloor c/64 \rfloor$ of the
  word of~$c$, and $a \ge 0$. Then \texttt{check\_bit} returns whether every third-stage prime
  admits the class of $(a, b, c)$, \lname{JPoints.stage3Keeps} (Definition~\ref{def:stage1Keeps}),
  and leaves the data of the third stage consistent for the row~$(a, b)$.
\end{lemma}

\begin{proof}
  \leanok
  \uses{def:refIsSquare, def:Tables.WF, def:CurveTables.WF, def:SiftData.WF}
  Unless the data are for the row, \texttt{row3\_setup} sets them, with the inverses of the tables
  of \texttt{init\_main}, which are those modulo~$p$ (Definitions~\ref{def:Tables.WF}
  and~\ref{def:refIsSquare}). The operands of the reductions stay below~$2^{16}$, where the
  reduction by the reciprocal is the remainder (\lname{JPoints.mod3\_eq}), and the entry of the row
  at the class is \lname{JPoints.CurveTables.admits3} of the class.
\end{proof}

\begin{lemma}[The contract of \texttt{check\_bits}]\label{lem:checkBits_spec}
  \lean{JPoints.checkBits_spec}
  \leanok
  \uses{def:firstIf, def:HasPointOver, def:stage1Keeps, def:colKeeps, def:CheckGlobals, def:Printf.emitPoint, def:IsStageSetup, def:checkOnePoint}
  Let the set-up of the stages hold, the data of the third stage be consistent, the residues be
  those of the word of index~$i$, $(a, b)$ a row of the box, and let the bits of the word
  \texttt{nums} lie within the height bound. Then \texttt{check\_bits} prints the points over the
  triples $(a, b, c)$ of the bits~$c$ of \texttt{nums} that pass \lname{JPoints.coprimeKeeps} and
  \lname{JPoints.stage3Keeps}, $c$ increasing. It keeps the data of the third stage consistent and
  changes no other field of the state but the counters and the output.
\end{lemma}

\begin{proof}
  \leanok
  \uses{def:orderedTriples, lem:mem_orderedTriples_iff, def:CheckGlobals.WF, def:checkOnePoint, lem:checkOnePoint_spec, lem:checkBit_spec, lem:relprime3_spec}
  The loop takes the bits of \texttt{nums} from the lowest, with the tests of
  Lemmas~\ref{lem:relprime3_spec} and~\ref{lem:checkBit_spec}. A bit that passes both gives a
  triple of the search (Lemma~\ref{lem:mem_orderedTriples_iff}), whose points the exact check
  prints (Lemma~\ref{lem:checkOnePoint_spec}).
\end{proof}

\begin{lemma}[The contract of \texttt{check\_point}]\label{lem:checkPoint_spec}
  \lean{JPoints.checkPoint_spec}
  \leanok
  \uses{def:firstIf, def:HasPointOver, def:stage1Keeps, def:colKeeps, def:CheckGlobals, def:Printf.emitPoint, def:IsStageSetup, def:sift}
  Let the set-up of the stages hold, the rows of the sieve tables be those of \texttt{row\_setup}
  for the row~$(a, b)$ of the box, the data of the third stage be consistent, the residues be
  those of the word of index~$i$ at the position~$k$, and the bits of \texttt{nums} lie within the
  height bound. Then \texttt{check\_point} prints the points over the triples $(a, b, c)$ of the
  bits~$c$ of \texttt{nums} that the second stage, the test of coprimality and the third stage
  admit (\lname{JPoints.stage2Keeps}, \lname{JPoints.coprimeKeeps}, \lname{JPoints.stage3Keeps}),
  $c$ increasing. It keeps the data of the third stage consistent and changes no other field of the
  state but the counters and the output.
\end{lemma}

\begin{proof}
  \leanok
  \uses{def:refIsSquare, def:CurveTables.WF, def:SieveTables.WF, lem:SieveTables.WF.getLsbD_eq_admits12, def:SiftData.WF, lem:checkBits_spec}
  For each second-stage prime~$p$ the word is and-ed with word $i \bmod p$ of the row of the class
  of~$(a, b)$, whose bit~$j$ is the class test of $(a, b, 64 i + j)$
  (Lemma~\ref{lem:SieveTables.WF.getLsbD_eq_admits12}, the test depending on the classes only); the
  loop stops early at the word~$0$, which no later prime changes. The bits left go to
  \texttt{check\_bits} (Lemma~\ref{lem:checkBits_spec}).
\end{proof}

\begin{lemma}[The contract of \texttt{scan\_words}]\label{lem:scanWords_spec}
  \lean{JPoints.scanWords_spec}
  \leanok
  \uses{def:firstIf, def:HasPointOver, def:stage1Keeps, def:colKeeps, def:CheckGlobals, def:Printf.emitPoint, def:IsStageSetup, def:sift}
  Let the precondition \lname{JPoints.ScanPre} of the second stage of each word hold: the set-up of
  the stages, the rows of the sieve tables and consistent data of the third stage for the
  row~$(a, b)$ of the box, the residues of the words of the positions and their bits within the height
  bound. Then \texttt{scan\_words} passes the nonzero words of the positions $[lo, hi)$ of the chunk
  to \texttt{check\_point}, in order, and prints the points that these find
  (\lname{JPoints.scanPoints}). It changes no other field of the state but the data of the third
  stage, the counters and the output.
\end{lemma}

\begin{proof}
  \leanok
  \uses{lem:checkPoint_spec}
  A loop over the positions: a word~$0$ has no points, each other one those of
  Lemma~\ref{lem:checkPoint_spec}, and under \texttt{-1} the loop stops at the first point.
\end{proof}

When the bit array holds exactly the bits that the first stage, the mask of the row, the masks of
the ends of the row and~$c_0$ leave, the points that \texttt{scan\_words} finds are those of the
range of words (\lname{JPoints.IsRowSetup.scanPoints\_eq\_rowRangePoints}).

\section{The first stage and the sieve of a chunk}\label{sec:stage1}

For a first-stage prime~$p$, word~$t$ of the bit array of a chunk, the word of index
$w_{\mathrm{low}} + t$ of the row for the first word~$w_{\mathrm{low}}$ of the chunk, is and-ed
with word $(w_{\mathrm{low}} + t) \bmod p$ of the row of~$(a, b)$ in the table of~$p$. The walk of~$p$ over the
chunk, which \texttt{init\_sift} computes once (\lname{JPoints.initSiftWalk}), splits the words
into a head from the word $\mathrm{first} \in [1, p]$ of the row with
$\mathrm{first} \equiv w_{\mathrm{low}} \bmod p$ up to the end of the period, $\mathrm{nper}$ full
periods of~$p$~words, and a tail; the macro \texttt{WALK} (\csource{j-sift.c}{914}{918}) runs
through them without a division. For a range of words of a chunk, \texttt{walk\_parts}
(\csource{j-sift.c}{895}{906}) makes the walk over the range from that of the chunk, with the
number of the full periods computed in doubles.

\formalnote{The model is that of the scalar walk, \texttt{VW = 1}, of the default build.}

\begin{lemma}[The walks of the first stage]\label{lem:walk_spec}
  \lean{JPoints.walk_spec, JPoints.walkParts_spec}
  \leanok
  \uses{def:IsStageSetup, def:Double.add, def:Double.ofDyadic}
  \begin{enumerate}
    \item Let~$w$ be a well-formed walk of the prime~$p$ (\lname{JPoints.WalkSpec.WF}:
      $1 \le \mathrm{first}$, $\mathrm{first} + \mathrm{head} \le p$, $\mathrm{tail} < p$, and
      $\mathrm{first} + \mathrm{head} = p$ when full periods or a tail follow), whose
      $\ell = \mathrm{head} + \mathrm{nper}\, p + \mathrm{tail}$ words from \texttt{surv0} on lie
      in the bit array. Then \texttt{WALK} with the operation~$\mathrm{op}$ sets word
      $\texttt{surv0} + t$, $t < \ell$, to $\mathrm{op}$ of its old value and of word
      $(\mathrm{first} + t) \bmod p$ of the row of the sieve table, and leaves the other words
      unchanged.
    \item Under the set-up of the stages (Definition~\ref{def:IsStageSetup}), for the~$n$-th
      sieving prime~$p$, $n < s_3$, a walk with $1 \le \mathrm{first} \le p$, a position
      $\mathit{off}$ and a length $\mathit{len} < 2^{48}$, both at most \texttt{array\_size},
      \texttt{walk\_parts} returns a well-formed walk over $\mathit{len}$ words whose first word is
      congruent to $\mathrm{first} + \mathit{off}$ modulo~$p$ (\lname{JPoints.WalkSpec.IsWalk}),
      its full periods counted exactly in doubles.
  \end{enumerate}
\end{lemma}

\begin{proof}
  \leanok
  \uses{def:SiftData.WF, def:Double.toRat, def:Double.ofDyadic, lem:Double.ofDyadic_exact, def:Double.add, def:Double.IsRounding, lem:Double.isRounding_ofDyadic, lem:Double.isRounding_add}
  (1) The head reads the words from~$\mathrm{first}$ on, each full period and the tail those from
  the start of the row on; as the head ends at the end of the period when more follows, word~$t$
  of the walk is word $(\mathrm{first} + t) \bmod p$. (2) The first word~$f$ is
  $\mathrm{first} + (\mathit{off} \bmod p)$, from the residues of the positions
  (Definition~\ref{def:SiftData.WF}), reduced into~$[1, p]$; the head has
  $h = \min(p - f, \mathit{len})$ words, and the tail the residue~$r$ modulo~$p$ of the number
  $\mathit{rem} = \mathit{len} - h$ of words after it. The full periods are
  $(\mathit{rem} - r) \cdot (1/p) + 0.5$ in doubles, truncated: the conversion is exact
  (Lemma~\ref{lem:Double.ofDyadic_exact}), and the roundings to nearest of the product and the sum
  (Lemmas~\ref{lem:Double.isRounding_ofDyadic} and~\ref{lem:Double.isRounding_add}) move the
  integer $(\mathit{rem} - r)/p$ by less than~$1/2$.
\end{proof}

\begin{lemma}[The contract of \texttt{sift0}]\label{lem:sift0_spec}
  \lean{JPoints.sift0_spec}
  \leanok
  \uses{def:firstIf, def:CheckGlobals, def:Printf.emitPoint, def:IsStageSetup, def:rowRangePoints, def:sift}
  Let the precondition \lname{JPoints.Sift0Pre} hold: what the sieve of a row reads
  (Definition~\ref{def:IsStageSetup}), a chunk~$k$ of the row, the rows of the sieve tables for the
  row~$(a, b)$ of the box, consistent data of the third stage, a bit array of
  $\texttt{array\_size} + 1$ words, and the mask and~$c_0$ of the row. Then \texttt{sift0} prints
  the points $\mathrm{rowRangePoints}(g, sd, a, b, w, w')$ of the words $[w, w')$ of the~$k$-th
  chunk (Definition~\ref{def:rowRangePoints}). It changes no other field of the state but the bit
  array, the data of the third stage, the counters and the output.
\end{lemma}

\begin{proof}
  \leanok
  \uses{def:orderedTriples, def:HasPointOver, def:stage1Keeps, def:colKeeps, def:CheckGlobals.WF, def:refIsSquare, def:CurveTables.WF, def:SieveTables.WF, lem:SieveTables.WF.getLsbD_eq_admits12, def:SiftData.WF, lem:walk_spec, lem:scanWords_spec}
  After the fill, the masks and the walks of the other first-stage primes over the chunk
  (Lemma~\ref{lem:walk_spec}, with the walks of \texttt{init\_sift}), bit~$j$ of the word of
  index~$i$ is set exactly when $c = 64 i + j$ passes the mask of the row, lies in the row, is
  not~$c_0$ and passes the first stage (Lemma~\ref{lem:SieveTables.WF.getLsbD_eq_admits12}). The
  scan of the words (Lemma~\ref{lem:scanWords_spec}) adds the second stage, the test of
  coprimality and the third, which together make \lname{JPoints.colKeeps}; and~$c_0$ is the only
  third coordinate of the row on the conic.
\end{proof}

\begin{lemma}[The contract of \texttt{sift}]\label{lem:sift_spec}
  \lean{JPoints.sift_spec}
  \leanok
  \uses{def:firstIf, def:rowPoints, def:CheckGlobals, def:Printf.emitPoint, def:SiftState.WF, def:IsSearchSetup, def:sift}
  If \lname{JPoints.IsSieveSetup} holds for the globals of the sieve by bands and the static data
  (Definition~\ref{def:IsSearchSetup}), the state is consistent
  (Definition~\ref{def:SiftState.WF}) and~$(a, b)$ is a row of the box, then \texttt{sift} prints
  the points $\mathrm{rowPoints}(a, b)$ of the row (Definition~\ref{def:rowPoints}) and keeps the
  state consistent.
\end{lemma}

\begin{proof}
  \leanok
  \uses{def:orderedTriples, def:HasPointOver, def:stage1Keeps, def:colKeeps, def:Cli.Input, def:Cli.Input.WF, def:CheckGlobals.WF, def:IsSearchGlobals, def:Tables.WF, def:SieveTables.WF, def:SiftData.WF, def:IsStageSetup, def:rowRangePoints, lem:rowSetup_spec, lem:sift0_spec}
  The set-up of the search gives that of a row (Definition~\ref{def:IsStageSetup}), with the
  globals of the exact check of an accepted command line (Definition~\ref{def:Cli.Input.WF}).
  \texttt{row\_setup} sets the rows (Lemma~\ref{lem:rowSetup_spec}), and the chunks in order print
  the points of their words (Lemma~\ref{lem:sift0_spec}), which together are those of the row;
  under \texttt{-1} the loop stops at the first point.
\end{proof}

\begin{lemma}[The contract of \texttt{sift\_range}]\label{lem:siftRange_spec}
  \lean{JPoints.siftRange_spec}
  \leanok
  \uses{def:firstIf, def:CheckGlobals, def:Printf.emitPoint, def:IsStageSetup, def:rowRangePoints, def:sift}
  Let the precondition \lname{JPoints.SiftRangePre} hold, that of \texttt{sift0} with~$lo$ among
  the words of the~$k$-th chunk and~$hi$ at most its end. Then \texttt{sift\_range} prints the
  points $\mathrm{rowRangePoints}(g, sd, a, b, lo, hi)$ of the words $[lo, hi)$. It changes no
  other field of the state but the bit array, the data of the third stage, the counters and the
  output.
\end{lemma}

\begin{proof}
  \leanok
  \uses{def:orderedTriples, def:HasPointOver, def:stage1Keeps, def:colKeeps, def:CheckGlobals.WF, def:refIsSquare, def:CurveTables.WF, def:SieveTables.WF, lem:SieveTables.WF.getLsbD_eq_admits12, def:SiftData.WF, lem:walk_spec, lem:scanWords_spec}
  As for Lemma~\ref{lem:sift0_spec}, with the fill and the masks of the positions of the range and
  the walks over the range that \texttt{walk\_parts} makes from those of the chunk
  (Lemma~\ref{lem:walk_spec}).
\end{proof}

\section{The rows of a first coordinate}\label{sec:rows}

A sieving prime $p \mid a$ modulo which the curve has no points at infinity admits no class
$(0, b, c)$ with $p \nmid b$ (Definition~\ref{def:refAdmits12}), so that the rows~$(a, b)$ with
$p \nmid b$ hold no points (Lemma~\ref{lem:rowKeeps_of_hasPointOver}). \texttt{row\_step}
(\csource{j-points.c}{1274}{1285}) returns the product~$m$ of these primes, or~$H + 1$ once the
product exceeds~$H$, and \texttt{run\_a} visits only the rows with $m \mid b$. In the box it sieves
by \texttt{sift} each of them that the condition at~$2$ leaves
(Lemma~\ref{lem:aliveKeeps_of_hasPointOver}); in the tube and the real region it sieves them by
bands (Definition~\ref{def:siftBands}). \texttt{run\_units} runs \texttt{run\_a} over the first
coordinates of the units of work: all $0 \le a \le H$, or for a monic quintic the squares, since
only those carry points (Lemma~\ref{lem:isSquare_of_mem_kummerPoints}).

\begin{definition}[The step of the rows]\label{def:rowStepProd}
  \lean{JPoints.rowStepPrimes, JPoints.rowStepProd}
  \leanok
  The primes of the step of the rows of the first coordinate~$a$, \lname{JPoints.rowStepPrimes},
  are the sieving primes~$p_n$, $n < s_3$, in the order of the plan, that divide~$a$ and modulo
  which the curve has no points at infinity (\texttt{has\_infinity}). \lname{JPoints.rowStepProd}
  is the product of those among the sieving primes of a list of indices: the product that
  \texttt{row\_step} has formed after these indices.
\end{definition}

\begin{lemma}[The contract of \texttt{row\_step}]\label{lem:run_rowStep}
  \lean{JPoints.run_rowStep, JPoints.rowStep_spec}
  \leanok
  \uses{def:rowStepProd, def:CheckGlobals}
  If $H \ge 1$ and the sieving primes are primes of the table, then \texttt{row\_step} returns
  for~$a$ the product~$m$ of the primes of the step of the rows of~$a$ when $m \le H$, and~$H + 1$
  otherwise; in particular $1 \le \texttt{row\_step}(a) \le H + 1$.
\end{lemma}

\begin{proof}
  \leanok
  The partial products grow, the primes of the table being at least~$3$, so that one of them
  exceeds~$H$ exactly when the whole product does.
\end{proof}

\begin{definition}[The rows of a first coordinate and the units of work]\label{def:runA}
  \lean{JPoints.runA, JPoints.runUnits}
  \leanok
  \uses{def:sift, def:siftBands, def:CheckGlobals}
  \lname{JPoints.runA} is \texttt{run\_a} (\csource{j-points.c}{1330}{1335}): for the first
  coordinate $a \ge 0$, the step $m = \texttt{row\_step}(a)$ and the rows~$(a, b)$ with $m \mid b$,
  from $b_0 = m$ for $a = 0$ and from $b_0 = -\lfloor H/m \rfloor\, m$ otherwise, up to~$H$; in the
  box each row sieved by \lname{JPoints.sift}, but for the rows of $a \ne 0$ whose class
  modulo~$64$ the condition at~$2$, when it is used, empties; in the tube or the real region the
  rows sieved by bands, \lname{JPoints.siftBands}; until a point is found under \texttt{-1}.
  \lname{JPoints.runUnits} is \texttt{run\_units} (\csource{j-points.c}{1385}{1386}):
  \texttt{run\_a} for the first coordinates of the units of work $k < \mathit{units}$ in order
  (\lname{JPoints.unitA}), until a point is found under \texttt{-1}, then the totals of the
  counters.
\end{definition}

\formalnote{The threads of \texttt{-t} are not modeled: the model runs the arm of one thread for
every number of threads, which prints the points in the order of the units, as \jpoints{} prints
them with threads; that the output is the same with threads is not proved, and the differential
tests compare runs with threads too.}

\begin{lemma}[The contract of \texttt{run\_a}]\label{lem:runA_spec}
  \lean{JPoints.runA_spec}
  \leanok
  \uses{def:firstIf, def:HasPointOver, def:mainTriples, def:CheckGlobals, def:Printf.emitPoint, def:SiftState.WF, def:IsSearchSetup, def:runA}
  If \lname{JPoints.IsSieveSetup} holds, the state is consistent and $0 \le a \le H$, then
  \texttt{run\_a} prints the points over the triples of the main search with the first
  coordinate~$a$ (\lname{JPoints.unitTriples}, Definition~\ref{def:mainTriples}), in their order,
  and keeps the state consistent.
\end{lemma}

\begin{proof}
  \leanok
  \uses{def:rowKeeps, def:rowPoints, lem:flatMap_rowPoints_rowsFrom, def:Cli.Input, def:Cli.Input.WF, def:CheckGlobals.WF, def:IsSearchGlobals, def:sift, lem:sift_spec, lem:run_rowStep, def:siftBands, lem:siftBands_spec}
  The rows visited are those with $m \mid b$ that the test of the rows keeps
  (Definition~\ref{def:rowKeeps}, with Lemma~\ref{lem:run_rowStep} for~$m$). A row sieved in the
  box prints its points (Lemma~\ref{lem:sift_spec}), a row sieved by bands those of
  Lemma~\ref{lem:siftBands_spec}, and the points of the rows in order are those over the triples of
  the first coordinate~$a$ (Lemma~\ref{lem:flatMap_rowPoints_rowsFrom}).
\end{proof}

\begin{lemma}[The contract of \texttt{run\_units}]\label{lem:runUnits_spec}
  \lean{JPoints.runUnits_spec}
  \leanok
  \uses{def:firstIf, def:HasPointOver, def:mainTriples, def:CheckGlobals, def:Printf.emitPoint, def:SiftState.WF, def:IsSearchSetup, def:runA}
  If \lname{JPoints.IsSieveSetup} holds, the state is consistent and the number of the units of
  work is $\lfloor \sqrt{H} \rfloor + 1$ for a monic quintic and~$H + 1$ otherwise, then
  \texttt{run\_units} prints the points over the triples of the main search off the conic
  (\lname{JPoints.mainTriples}), in their order.
\end{lemma}

\begin{proof}
  \leanok
  \uses{lem:flatMap_range_unitTriples_unitA, def:CheckGlobals.WF, def:IsSearchGlobals, lem:runA_spec}
  The units are the first coordinates $0 \le a \le H$, or the squares $a = k^2$ with
  $k \le \lfloor \sqrt{H} \rfloor$ for a monic quintic; the points over the triples of the units
  in order are those of the main search (Lemma~\ref{lem:flatMap_range_unitTriples_unitA}), and each
  unit prints its own (Lemma~\ref{lem:runA_spec}).
\end{proof}

\section{The double points}\label{sec:doublePoints}

The triples of the conic $b^2 = 4 a c$ other than $(0, 0, 1)$ are those of the double points
$x = t/s$, $(s^2, 2 s t, t^2)$ with $s \ge 1$ and $\gcd(s, t) = 1$ (\lname{JPoints.doubleTriples}).
\texttt{find\_double\_points} (\csource{j-points.c}{1192}{1197}) searches them row by row: for each
first coordinate~$s$ with $s^2 \le H$, the second coordinates
$|t| \le \min(H/(2 s), \lfloor \sqrt{H} \rfloor)$, with a sieve of the pairs at the sieving primes.
For a curve that is not a monic quintic it skips every~$s$ that a sieving prime modulo which the
curve has no points at infinity divides; for a monic quintic only the squares~$s$ are first
coordinates (\lname{JPoints.isSquare\_of\_isSquare\_sexticForm}). The sieve, the skip and the
squares are sound because a point over the triple makes $F(t, s) = s^6 f(t/s)$ a square.

\begin{lemma}[The points over the double points]\label{lem:isSquare_sexticForm_of_hasPointOver}
  \lean{JPoints.isSquare_sexticForm_of_hasPointOver}
  \leanok
  \uses{def:auxA2, def:HasPointOver}
  If $s \ne 0$ and there is a point over the triple $(s^2, 2 s t, t^2)$
  (Definition~\ref{def:HasPointOver}), then $F(t, s) = s^6 f(t/s)$ is a square in~$\ZZ$.
\end{lemma}

\begin{proof}
  \leanok
  \uses{def:kummer, def:Lifts, def:kummerRoots, lem:mem_kummerRoots_iff}
  The fourth coordinate of the point is a root~$d/m$ of the Kummer equation over the triple
  (Lemma~\ref{lem:mem_kummerRoots_iff}). Over $(m s^2, 2 m s t, m t^2, d)$, the identity
  $m^4 s^4 F(t, s) = t^2 A_2 + s t\, C + s^2 B_2$
  (\lname{JPoints.pow\_four\_mul\_pow\_four\_mul\_sexticForm\_eq}) and the lifting test, which makes
  $A_2 = \alpha^2$, $B_2 = \beta^2$ and $C = 2 \alpha \beta$ (Definition~\ref{def:Lifts}), give
  $m^4 s^4 F(t, s) = (t \alpha + s \beta)^2$, so that $F(t, s)$ is a square.
\end{proof}

\begin{definition}[The sieve of the double points]\label{def:DoubleSetup}
  \lean{JPoints.DoubleSetup, JPoints.Stype.passes}
  \leanok
  \uses{def:CheckGlobals, def:CheckGlobals.WF, def:refIsSquare, def:Tables.WF}
  What \texttt{find\_double\_points} reads of the globals of the sieve,
  \lname{JPoints.DoubleSetup}: the globals of the exact check as \texttt{main} sets them
  (Definition~\ref{def:CheckGlobals.WF}), the tables of \texttt{init\_main}
  (Definition~\ref{def:Tables.WF}), at most \texttt{NUM\_PRIMES} sieving primes of the table, and
  at each of them the points at infinity and the values of~$f$ that are squares as the references
  say (Definition~\ref{def:refIsSquare}). The record of the sieve for a sieving prime~$p$ and the
  first coordinate~$s$ (\lname{JPoints.doubleRecord}) holds the inverse~$s'$ of $s \bmod p$,
  or~$0$ when $p \mid s$, and the table that decides: \texttt{has\_infinity} at~$p$ when
  $p \mid s$, the row of~$p$ of \texttt{is\_f\_square} otherwise. It passes the second
  coordinate~$t$, \lname{JPoints.Stype.passes}, when its table has the entry
  $s' (t \bmod p) \bmod p$.
\end{definition}

\begin{lemma}[Soundness of the sieve of the double points]\label{lem:DoubleSetup.passes_of_points_ne_nil}
  \lean{JPoints.DoubleSetup.passes_of_points_ne_nil}
  \leanok
  \uses{def:HasPointOver, def:CheckGlobals, def:DoubleSetup}
  If \lname{JPoints.DoubleSetup} holds, $a \ne 0$, $\gcd(a, b) = 1$ and the exact check finds a
  point over $(a^2, 2 a b, b^2)$, then at every sieving prime the record of~$a$ passes~$b$.
\end{lemma}

\begin{proof}
  \leanok
  \uses{def:kummer, def:auxA2, def:Lifts, def:IsNormalized, def:InBounds, def:IsPoint, def:kummerRoots, def:refIsSquare, def:Tables.WF, lem:isSquare_sexticForm_of_hasPointOver}
  $F(b, a)$ is a square (Lemma~\ref{lem:isSquare_sexticForm_of_hasPointOver}). For $p \nmid a$,
  $F(b, a) = a^6 f(b/a)$ makes $f(b a^{-1})$ a square modulo~$p$; for $p \mid a$,
  $F(b, a) \equiv f_6 b^6 \bmod p$ with $p \nmid b$ makes~$f_6$ a square modulo~$p$, so that the
  curve has points at infinity modulo~$p$ (Definition~\ref{def:refIsSquare}). The same argument
  shows that \texttt{find\_double\_points} skips no first coordinate with a point
  (\lname{JPoints.DoubleSetup.hasInfinity\_of\_points\_ne\_nil}).
\end{proof}

\begin{lemma}[The contract of \texttt{find\_double\_points}]\label{lem:findDoublePoints_spec_of_doubleSetup}
  \lean{JPoints.findDoublePoints_spec_of_doubleSetup, JPoints.findDoublePoints_spec}
  \leanok
  \uses{def:firstIf, def:HasPointOver, def:mainTriples, def:CheckGlobals, def:Printf.emitPoint, def:IsSearchSetup, def:DoubleSetup, def:Double.add, def:Double.ofDyadic, def:checkOnePoint}
  If \lname{JPoints.DoubleSetup} holds for the globals of the sieve, which it does under
  \lname{JPoints.IsSieveSetup} (Definition~\ref{def:IsSearchSetup}), then
  \texttt{find\_double\_points} prints the points over the triples of the double points
  (\lname{JPoints.doubleTriples}, Definition~\ref{def:mainTriples}), in their order.
\end{lemma}

\begin{proof}
  \leanok
  \uses{def:poly, def:kummer, def:auxA2, def:Lifts, def:IsNormalized, def:InBounds, def:IsPoint, def:orderedTriples, def:kummerRoots, lem:mem_orderedTriples_iff, def:Cli.Input, def:Cli.Input.WF, def:CheckGlobals.WF, def:checkOnePoint, lem:checkOnePoint_spec, def:IsSearchGlobals, def:refIsSquare, def:Tables.WF, def:CurveTables.WF, lem:isSquare_sexticForm_of_hasPointOver, lem:DoubleSetup.passes_of_points_ne_nil, def:Double.toRat, def:Double.ofDyadic, lem:Double.ofDyadic_exact, def:Double.add}
  The bound $\lfloor \sqrt{H} \rfloor$, which \texttt{find\_double\_points} computes as
  \texttt{floor(sqrt(height))} in doubles, is exact for $H \le \texttt{MAX\_HEIGHT}$, $H$
  converting exactly (Lemma~\ref{lem:Double.ofDyadic_exact}). The loop over the first coordinates
  runs over the squares for a monic quintic, whose other rows have no points, and over all
  $1 \le s \le \lfloor \sqrt{H} \rfloor$ otherwise, skipping only rows without points; in each row
  the coprime pairs that pass the sieve (Lemma~\ref{lem:DoubleSetup.passes_of_points_ne_nil}) go to
  the exact check (Lemma~\ref{lem:checkOnePoint_spec}), their triples being those of the search
  (Lemma~\ref{lem:mem_orderedTriples_iff}).
\end{proof}
